We examine a scenario where a computation \( f(x): X \rightarrow Y \) is outsourced between a client \( \mathcal{C} \), who possesses the input \( x \), and a server \( \mathcal{S} \). The ownership of the function \( f(x): X \rightarrow Y \) varies: it may reside with the client \( \mathcal{C} \), in line with client-focused frameworks like Infrastructure as a Service (IaaS) \cite{serrano2015infrastructure} or Platform as a Service (PaaS) \cite{keller2010platform}, or it could belong to the server \( \mathcal{S} \), as seen in services like Machine Learning as a Service (MLaaS) \cite{ribeiro2015mlaas}, where the server is responsible for providing the model data (e.g., the weights in a Deep Neural Network). In our model, we assume that if part of the data \( X \) belongs to the server \( \mathcal{S} \), it is known to both parties. However, data originating from the client \( \mathcal{C} \) must be kept confidential. The primary computation we focus on is matrix operations like matrix multiplication and convolution. We remark that extending DataSeal to other computations would require new ABFT techniques.

The threat model we advance considers the possibility of the server \( \mathcal{S} \) being malevolent or compromised, possessing the ability to alter ciphertexts. This extends beyond the conventional semi-honest model, which perceives \( \mathcal{S} \) as simply honest but curious (HBC) \cite{le2023privacy,dong2022fusion, badolato2022privacy}. While the HBC approach offers some promise, it is vulnerable to the threat of malign servers. Our model is designed to address proactive adversarial actions, including data manipulation and unauthorized data access. In this context, the adversarial server's objectives are twofold: (i) Violating data privacy by accessing and potentially understanding the client's data $x$; (ii) Compromising computation integrity by modifying the data in a way that changes its meaning or accuracy, all while avoiding detection. 

To ensure the security of DataSeal under this threat model, we define a formal security game that captures the server's ability to forge computation results. DataSeal's security is based on the unforgeability of the encoded matrices and checksums transmitted to the server. The adversary’s goal is to produce a forged result that passes client-side verification without access to the secret key or random scalar. As detailed in Appendix A, the probability that an adversary can successfully generate a valid forged result is shown to be negligible under standard cryptographic assumptions. Please refer to Appendix A for a detailed description of the security guarantees provided by DataSeal.



% \subsection{Objectives}
% Our objectives are: (1) The client wants to protect the confidentiality. (2) The client wants to ensure that the computation is correct through verification.